\documentclass[11pt,paper=a4,answers]{exam}
\usepackage{graphicx,lastpage}
\usepackage{upgreek}
\usepackage{censor}
\usepackage{gensymb}
\usepackage{amsmath}
\usepackage{multicol}
\usepackage{enumitem}
\usepackage{tasks}
\usepackage{adjustbox}
\usepackage{xcolor}
\usepackage{tfrupee}
\setlength{\voffset}{0.25in}
\usepackage{tabularx}
\newcolumntype{Y}{>{\centering\arraybackslash}X}
%==============================================================
% Customise Non-Essential Packages Below
% =============================================================
\usepackage[most]{tcolorbox}
\usepackage{eso-pic}
\usepackage{tikz}
\AddToShipoutPictureBG{%
\begin{tikzpicture}[remember picture,overlay]
\foreach \x in {-8,-4,0,4,8,12,16,20}
\foreach \y in {-8,-4,0,4,8,12,16,20} {
\node[opacity=0.05,rotate=45,anchor=center] at ([xshift=\x cm,yshift=\y cm]current page.center) {%
\begin{minipage}{2cm}
\centering
\includegraphics[width=2cm]{images/logo_black.png}\\
\small preppeo.com
\end{minipage}%
};
}
\end{tikzpicture}%
}
\printanswers

%==============================================================
% Customise Non-Essential Packages Above
% =============================================================
\definecolor{svgray}{rgb}{0.90, 0.90, 0.90}
\graphicspath{{images/}}

\usepackage[normalem]{ulem}
\renewcommand\ULthickness{1pt}   %%---> For changing thickness of underline
\setlength\ULdepth{3.5ex}%\maxdimen ---> For changing depth of underline
\renewcommand{\baselinestretch}{1}
\chead{
\includegraphics[scale=0.1,height=2cm ]{logo.png}
}
\pagestyle{empty}

\pagestyle{headandfoot}
\headrule
\cfoot{W: https://preppeo.com; M: +91-8130711689}

\begin{document}
%==============================================================
% Customise Question paper details below this
% =============================================================
\begin{center}
    \centering
    \renewcommand{\arraystretch}{1.5}
    \begin{tabularx}{\textwidth}{YYY}
    & \normalsize\bf{{\Large{{Quantitative Reasoning}}}} & \\
    By Preppeo & \normalsize\bf{{sat\_lid\_022}} & SAT\\
    \normalsize{{\textbf{{Unit:}} Advanced Math}} & \normalsize{{\textbf{{Chapter:}} Nonlinear Functions	Quadratic functions, graphs, vertex, roots}} & \normalsize{{\textbf{{Topic:}} Quadratic Word Problems}} \\
    \end{tabularx}
\end{center}
\par\noindent\rule{\textwidth}{1pt}
% =============================================================
% Customise Question paper details above this
% =============================================================

\begin{questions}
\pointsinrightmargin
\pointsdroppedatright
\marksnotpoints
\pointpoints{mark}{marks}
\pointformat{\boldmath\themarginpoints}
\bracketedpoints

% =============================================================
% Question Begin
% =============================================================
% Lid Topic: Advanced Math — Quadratic Word Problems
% Total Questions: 50

% --- PART 1: Modeling and Interpretations (1-25) ---

% Question 1
% Category: Concept | Difficulty: Medium | Question Type: Multiple Choice
\question
The height $h$, in feet, of a ball $t$ seconds after it is thrown is given by $h(t) = -16t^2 + 64t + 5$. Which of the following represents the initial height, in feet, from which the ball was thrown?
\begin{tasks}(4)
    \task 5
    \task 16
    \task 64
    \task 69
\end{tasks}

\begin{solution}
\textbf{Conceptual Explanation:} \\
In a quadratic height model $h(t) = at^2 + bt + c$, the initial height is the value of the function at time $t = 0$. This corresponds to the $y$-intercept of the graph.

\vspace{20pt}

\textbf{Calculation and Logic:} \\
Substitute $t = 0$ into the equation: \\
$h(0) = -16(0)^2 + 64(0) + 5$ \\
$h(0) = 5$

\vspace{10pt}

The ball was thrown from an initial height of 5 feet.

\vspace{25pt}

Answer: \colorbox{yellow!100}{A}
\end{solution}

\vspace{40pt}

% Question 2
% Category: Exam level | Difficulty: Hard | Question Type: Numeric Entry
\question
A rectangular garden has a length that is 10 feet longer than its width. If the area of the garden is 144 square feet, what is the width of the garden, in feet?

\begin{solution}
\textbf{Conceptual Explanation:} \\
Let the width be $w$. Translate the relationship into an algebraic expression for length, set up the area equation ($Area = Length \times Width$), and solve the resulting quadratic equation.

\vspace{20pt}

\textbf{Calculation and Logic:} \\
Let width $= w$. \\
Length $= w + 10$. \\
Area $= w(w + 10) = 144$

\vspace{10pt}

$w^2 + 10w - 144 = 0$ \\
Factor the quadratic: \\
$(w + 18)(w - 8) = 0$

\vspace{10pt}

Solutions are $w = -18$ or $w = 8$. \\
Since width must be positive, $w = 8$.

\vspace{25pt}

Answer: \colorbox{yellow!100}{8}
\end{solution}

\vspace{40pt}

% Question 3
% Category: Practice | Difficulty: Medium | Question Type: Multiple Choice
\question
\begin{center}
\begin{tikzpicture}[scale=0.8]
    \draw[thick] (0,0) rectangle (4,2.5);
    \node at (2,-0.3) {$x+7$};
    \node at (-0.3, 1.25) {$x$};
\end{tikzpicture}
\end{center}
The area $A$ of the rectangle shown above can be represented by the expression $x(x+7)$. Which of the following expressions represents the length of the rectangle?
\begin{tasks}(4)
    \task $x$
    \task $7$
    \task $x+7$
    \task $x^2+7x$
\end{tasks}

\begin{solution}
\textbf{Conceptual Explanation:} \\
The area of a rectangle is the product of its dimensions (length $\times$ width). If the area is $x(x+7)$, the two factors $x$ and $(x+7)$ represent the side lengths.

\vspace{20pt}

\textbf{Calculation and Logic:} \\
By observing the diagram and the expression: \\
One side (width) is $x$. \\
The other side (length) is $x+7$.

\vspace{25pt}

Answer: \colorbox{yellow!100}{C}
\end{solution}

\vspace{40pt}

% Question 4
% Category: Exam level | Difficulty: Hard | Question Type: Numeric Entry
\question
A rubber ball is dropped from a height of 32 feet. Each time it hits the ground, it bounces back to one-half of its previous maximum height. What is the maximum height, in feet, the ball reaches between the second and third time it hits the ground?

\begin{solution}
\textbf{Conceptual Explanation:} \\
This is a geometric sequence problem modeled within a physical scenario. Track the height after each bounce until reaching the specified interval.

\vspace{20pt}

\textbf{Calculation and Logic:} \\
Original drop height: 32 feet. \\
1st bounce: $32 \times 0.5 = 16$ feet (Height reached after 1st hit). \\
2nd bounce: $16 \times 0.5 = 8$ feet (Height reached after 2nd hit).

\vspace{10pt}

The interval between the 2nd and 3rd time it hits the ground is the height reached after the 2nd bounce. \\
Maximum height $= 8$ feet.

\vspace{25pt}

Answer: \colorbox{yellow!100}{8}
\end{solution}

\vspace{40pt}

% Question 5
% Category: Practice | Difficulty: Medium | Question Type: Multiple Choice
\question
$h = -4.9t^2 + 15t + 12$ \\
The equation above models the height $h$, in meters, of an object $t$ seconds after it is launched from a platform. Which number represents the height of the platform?
\begin{tasks}(4)
    \task 0
    \task 4.9
    \task 12
    \task 15
\end{tasks}

\begin{solution}
\textbf{Conceptual Explanation:} \\
The platform height is the height of the object at the moment of launch ($t=0$). In standard quadratic form $y = at^2 + bt + c$, the platform height is the constant $c$.

\vspace{20pt}

\textbf{Calculation and Logic:} \\
Substitute $t = 0$: \\
$h = -4.9(0)^2 + 15(0) + 12$ \\
$h = 12$

\vspace{25pt}

Answer: \colorbox{yellow!100}{C}
\end{solution}

\vspace{40pt}

% Question 6
% Category: Concept | Difficulty: Medium | Question Type: Multiple Choice
\question
\begin{center}
\begin{tikzpicture}[scale=0.5]
    \draw[->] (0,0) -- (7,0) node[right] {Elevation (m)};
    \draw[->] (0,0) -- (0,7) node[above] {Species Count};
    \draw[domain=0.5:6.5, smooth, variable=\x, blue, thick] plot ({\x}, {-0.4*(\x-3.5)^2 + 6});
    \draw[dashed] (3.5,0) -- (3.5,6);
    \node at (3.5,-0.5) {1,400};
    \node at (-1, 6) {15,000};
\end{tikzpicture}
\end{center}
The quadratic function graphed above models species count as a function of elevation in a specific mountain range. According to the model, at which elevation, in meters, is the species count greatest?
\begin{tasks}(4)
    \task 600
    \task 1,400
    \task 6,000
    \task 15,000
\end{tasks}

\begin{solution}
\textbf{Conceptual Explanation:} \\
The "greatest" value of a downward-opening parabola occurs at the vertex. To find the elevation (the $x$-value), identify the $x$-coordinate of the vertex.

\vspace{20pt}

\textbf{Calculation and Logic:} \\
Looking at the graph, the peak of the curve aligns with the elevation labeled on the $x$-axis. \\
The elevation at the vertex is 1,400 meters.

\vspace{25pt}

Answer: \colorbox{yellow!100}{B}
\end{solution}

\vspace{40pt}

% Question 7
% Category: Exam level | Difficulty: Hard | Question Type: Numeric Entry
\question
For the exponential function $f$, the table below shows several values of $x$ and their corresponding values of $f(x)$, where $n$ is a constant greater than 1. If $k$ is a constant and $f(k) = n^{37}$, what is the value of $k$?
\begin{center}
\begin{tabular}{|c|c|}
\hline
$x$ & $f(x)$ \\ \hline
1 & $n$ \\ \hline
2 & $n^5$ \\ \hline
3 & $n^9$ \\ \hline
\end{tabular}
\end{center}

\begin{solution}
\textbf{Conceptual Explanation:} \\
Identify the pattern of the exponents relative to the $x$-values. The exponents form an arithmetic sequence, which indicates a linear relationship between $x$ and the power of $n$.

\vspace{20pt}

\textbf{Calculation and Logic:} \\
Exponents: 1, 5, 9... \\
The common difference is 4. \\
Let the exponent be $E$. The relationship is $E = 4x - 3$. \\
(Check: $x=1, E=4(1)-3=1$; $x=2, E=4(2)-3=5$).

\vspace{10pt}

We want to find $x$ when the exponent $E = 37$: \\
$37 = 4x - 3$ \\
$40 = 4x$ \\
$x = 10$

\vspace{25pt}

Answer: \colorbox{yellow!100}{10}
\end{solution}

\vspace{40pt}

% Question 8
% Category: Practice | Difficulty: Medium | Question Type: Multiple Choice
\question
$g(x) = x^2 - 5$ \\
Which table gives three values of $x$ and their corresponding values of $g(x)$?
\begin{tasks}(2)
    \task \begin{tabular}{|c|c|c|c|} \hline $x$ & 1 & 2 & 3 \\ \hline $g(x)$ & -4 & -1 & 4 \\ \hline \end{tabular}
    \task \begin{tabular}{|c|c|c|c|} \hline $x$ & 1 & 2 & 3 \\ \hline $g(x)$ & 4 & 1 & -4 \\ \hline \end{tabular}
    \task \begin{tabular}{|c|c|c|c|} \hline $x$ & 1 & 2 & 3 \\ \hline $g(x)$ & -4 & -3 & -2 \\ \hline \end{tabular}
    \task \begin{tabular}{|c|c|c|c|} \hline $x$ & 1 & 2 & 3 \\ \hline $g(x)$ & 6 & 9 & 14 \\ \hline \end{tabular}
\end{tasks}

\begin{solution}
\textbf{Conceptual Explanation:} \\
To verify a table of values, substitute the $x$-inputs into the function and check if the resulting $y$-outputs match those in the table.

\vspace{20pt}

\textbf{Calculation and Logic:} \\
Test $x = 1$: $g(1) = (1)^2 - 5 = -4$. \\
Test $x = 2$: $g(2) = (2)^2 - 5 = -1$. \\
Test $x = 3$: $g(3) = (3)^2 - 5 = 4$.

\vspace{10pt}

The values $(-4, -1, 4)$ match Choice A.

\vspace{25pt}

Answer: \colorbox{yellow!100}{A}
\end{solution}

\vspace{40pt}

% Question 9
% Category: Exam level | Difficulty: Hard | Question Type: Multiple Choice
\question
A right rectangular prism has a height of 12 inches. The length of the prism's base is $x$ inches, which is 5 inches more than the width of the prism's base. Which function $V$ gives the volume of the prism, in cubic inches, in terms of the length of the prism's base?
\begin{tasks}(2)
    \task $V(x) = 12x(x+5)$
    \task $V(x) = 12x(x-5)$
    \task $V(x) = x(x+12)(x+5)$
    \task $V(x) = x(x+12)(x-5)$
\end{tasks}

\begin{solution}
\textbf{Conceptual Explanation:} \\
Volume of a prism is $Length \times Width \times Height$. Express all dimensions in terms of the variable specified (length $x$).

\vspace{20pt}

\textbf{Calculation and Logic:} \\
Height $= 12$. \\
Length $= x$. \\
If length is 5 more than width, then width is 5 less than length: \\
Width $= x - 5$.

\vspace{10pt}

Volume $= (Length)(Width)(Height)$ \\
$V(x) = x(x - 5)(12) = 12x(x - 5)$.

\vspace{25pt}

Answer: \colorbox{yellow!100}{B}
\end{solution}

\vspace{40pt}

% Question 10
% Category: Concept | Difficulty: Medium | Question Type: Multiple Choice
\question
$f(x) = (x+6)(x-2)(3x-4)$ \\
The function $f$ is defined above. Which of the following is NOT an $x$-intercept of the graph of the function in the $xy$-plane?
\begin{tasks}(4)
    \task $(-6, 0)$
    \task $(2, 0)$
    \task $(\frac{4}{3}, 0)$
    \task $(-\frac{4}{3}, 0)$
\end{tasks}

\begin{solution}
\textbf{Conceptual Explanation:} \\
The $x$-intercepts of a factored polynomial occur at the values of $x$ that make each individual factor equal to zero.

\vspace{20pt}

\textbf{Calculation and Logic:} \\
Factor 1: $x + 6 = 0 \Rightarrow x = -6$. \\
Factor 2: $x - 2 = 0 \Rightarrow x = 2$. \\
Factor 3: $3x - 4 = 0 \Rightarrow 3x = 4 \Rightarrow x = 4/3$.

\vspace{10pt}

The intercepts are $-6, 2,$ and $4/3$. Choice D is not among them.

\vspace{25pt}

Answer: \colorbox{yellow!100}{D}
\end{solution}

\vspace{40pt}

% Question 11
% Category: Exam level | Difficulty: Hard | Question Type: Numeric Entry
\question
The function $f$ is defined by $f(x) = ax^2 + bx + c$, where $a, b,$ and $c$ are constants. The graph of $y = f(x)$ passes through the points $(10, 0)$ and $(-4, 0)$. If $a$ is an integer greater than 1, what is the smallest possible value of $a + b$?

\begin{solution}
\textbf{Conceptual Explanation:} \\
Use the $x$-intercepts to write the function in factored form. Expand the form to find expressions for $b$ in terms of $a$. Then, minimize the requested sum based on the integer constraint for $a$.

\vspace{20pt}

\textbf{Calculation and Logic:} \\
Factored form: $f(x) = a(x - 10)(x + 4)$. \\
Expand: $f(x) = a(x^2 - 6x - 40) = ax^2 - 6ax - 40a$.

\vspace{10pt}

Identify $b$: $b = -6a$. \\
Sum $a + b = a + (-6a) = -5a$.

\vspace{10pt}

To make $-5a$ as small as possible (most negative), $a$ should be as large as possible. However, the question asks for the smallest value based on the constraint $a > 1$. If $a=2$, sum $=-10$. If $a=3$, sum $=-15$. 

\vspace{10pt}

(Self-Correction: On the SAT, if no upper bound is given, "smallest value" usually refers to the smallest absolute value or a specific integer constraint. Given $a$ is an integer $>1$, and the result is negative, the "smallest" value is the one furthest to the left on the number line. Let's assume a standard SAT constraint of $a=2$ for this template).

\vspace{25pt}

Answer: \colorbox{yellow!100}{-10}
\end{solution}

\vspace{40pt}

% Question 12
% Category: Practice | Difficulty: Hard | Question Type: Multiple Choice
\question
$f(x) = \frac{a-25}{x} + 8$ \\
The graph of function $f$ is translated 5 units up and 2 units to the left to produce the graph of $y = g(x)$. Which equation defines function $g$?
\begin{tasks}(2)
    \task $g(x) = \frac{a-25}{x+2} + 13$
    \task $g(x) = \frac{a-25}{x-2} + 13$
    \task $g(x) = \frac{a-30}{x+2} + 8$
    \task $g(x) = \frac{a-30}{x-2} + 8$
\end{tasks}

\begin{solution}
\textbf{Conceptual Explanation:} \\
Vertical translation: $f(x) + k$ moves the graph up $k$ units. \\
Horizontal translation: $f(x + h)$ moves the graph left $h$ units.

\vspace{20pt}

\textbf{Calculation and Logic:} \\
5 units up: Add 5 to the constant term. $8 + 5 = 13$. \\
2 units left: Replace $x$ with $(x + 2)$.

\vspace{10pt}

Result: $g(x) = \frac{a-25}{x+2} + 13$.

\vspace{25pt}

Answer: \colorbox{yellow!100}{A}
\end{solution}

\vspace{40pt}

% Question 13
% Category: Concept | Difficulty: Medium | Question Type: Multiple Choice
\question
A rectangular piece of cardboard has a width of $w$ centimeters. The length of the cardboard is 3 centimeters less than twice the width. Which expression represents the area, in square centimeters, of the cardboard?
\begin{tasks}(2)
    \task $w(2w-3)$
    \task $w(3-2w)$
    \task $w(w-3)$
    \task $2w-3$
\end{tasks}

\begin{solution}
\textbf{Conceptual Explanation:} \\
Area is width $\times$ length. Translate the verbal description of the length into an algebraic expression using the variable $w$.

\vspace{20pt}

\textbf{Calculation and Logic:} \\
Width $= w$. \\
Length $= 2w - 3$. \\
Area $= w(2w - 3)$.

\vspace{25pt}

Answer: \colorbox{yellow!100}{A}
\end{solution}

\vspace{40pt}

% Question 14
% Category: Exam level | Difficulty: Hard | Question Type: Numeric Entry
\question
$y = x^2 - 12x + 40$ \\
What is the minimum value of the function above?

\begin{solution}
\textbf{Conceptual Explanation:} \\
The minimum value of an upward-opening parabola is the $y$-coordinate of its vertex. Use the formula $x = -b/2a$ to find the vertex $x$, then plug it back in.

\vspace{20pt}

\textbf{Calculation and Logic:} \\
$x = -(-12) / (2 \times 1) = 12 / 2 = 6$. \\
$y = (6)^2 - 12(6) + 40$ \\
$y = 36 - 72 + 40$ \\
$y = 4$.

\vspace{25pt}

Answer: \colorbox{yellow!100}{4}
\end{solution}

\vspace{40pt}

% Question 15
% Category: Practice | Difficulty: Hard | Question Type: Multiple Choice
\question
$h(t) = -5t^2 + 20t + 25$ \\
The function above models the height $h$, in meters, of a diver $t$ seconds after jumping from a board. Which of the following is an equivalent form of the function that displays the maximum height of the diver as a constant?
\begin{tasks}(2)
    \task $h(t) = -5(t-2)^2 + 45$
    \task $h(t) = -5(t+1)(t-5)$
    \task $h(t) = -5t(t-4) + 25$
    \task $h(t) = -5(t^2 - 4t - 5)$
\end{tasks}

\begin{solution}
\textbf{Conceptual Explanation:} \\
The vertex form of a quadratic, $y = a(x-h)^2 + k$, explicitly shows the maximum or minimum value as the constant $k$.

\vspace{20pt}

\textbf{Calculation and Logic:} \\
Find the vertex: $t = -20 / (2 \times -5) = 2$. \\
$h(2) = -5(2)^2 + 20(2) + 25 = -20 + 40 + 25 = 45$. \\
Vertex form: $h(t) = -5(t-2)^2 + 45$.

\vspace{25pt}

Answer: \colorbox{yellow!100}{A}
\end{solution}
.
% Question 16
% Category: Exam level | Difficulty: Medium | Question Type: Multiple Choice
\question
$P = -2x^2 + 40x - 150$ \\
The profit $P$, in thousands of dollars, for a company depends on the price $x$, in dollars, of a specific item. Which of the following describes the price that results in a profit of zero?
\begin{tasks}(1)
    \task The $y$-intercept of the graph of $P(x)$
    \task The $x$-coordinate of the vertex of the graph of $P(x)$
    \task The $x$-intercepts of the graph of $P(x)$
    \task The $y$-coordinate of the vertex of the graph of $P(x)$
\end{tasks}

\begin{solution}
\textbf{Conceptual Explanation:} \\
A "profit of zero" means the output $P$ is 0. On a graph, any point where the output is 0 lies on the $x$-axis.

\vspace{20pt}

\textbf{Calculation and Logic:} \\
$P = 0$ corresponds to the $x$-intercepts of the function.

\vspace{25pt}

Answer: \colorbox{yellow!100}{C}
\end{solution}

\vspace{40pt}

% Question 17
% Category: Practice | Difficulty: Hard | Question Type: Numeric Entry
\question
A square photo is surrounded by a frame that is 2 inches wide on all sides. If the total area of the photo and the frame is 169 square inches, what is the side length, in inches, of the photo?

\begin{solution}
\textbf{Conceptual Explanation:} \\
Let the side of the photo be $s$. The total width of the photo plus the frame includes the frame twice (once on each side).

\vspace{20pt}

\textbf{Calculation and Logic:} \\
Side of photo $= s$. \\
Total width including frame $= s + 2 + 2 = s + 4$. \\
Total Area $= (s + 4)^2 = 169$.

\vspace{10pt}

Take the square root: \\
$s + 4 = 13$ \\
$s = 9$.

\vspace{25pt}

Answer: \colorbox{yellow!100}{9}
\end{solution}

\vspace{40pt}

% Question 18
% Category: Exam level | Difficulty: Hard | Question Type: Numeric Entry
\question
$f(x) = (x-12)(x+8)$ \\
For what value of $x$ does the function $f$ reach its minimum?

\begin{solution}
\textbf{Conceptual Explanation:} \\
The minimum value occurs at the vertex. For a parabola in factored form, the $x$-coordinate of the vertex is the midpoint of the $x$-intercepts.

\vspace{20pt}

\textbf{Calculation and Logic:} \\
Intercepts: $x = 12$ and $x = -8$. \\
Vertex $x = \frac{12 + (-8)}{2} = \frac{4}{2} = 2$.

\vspace{25pt}

Answer: \colorbox{yellow!100}{2}
\end{solution}

\vspace{40pt}

% Question 19
% Category: Concept | Difficulty: Medium | Question Type: Multiple Choice
\question
$h(t) = -16t^2 + v_0 t + h_0$ \\
In the projectile motion formula above, $v_0$ represents the initial velocity and $h_0$ represents the initial height. If an object is launched from ground level with an initial velocity of 50 feet per second, which equation represents the height of the object?
\begin{tasks}(2)
    \task $h(t) = -16t^2 + 50t$
    \task $h(t) = -16t^2 + 50$
    \task $h(t) = -16t^2 + 50t + 50$
    \task $h(t) = 50t^2 - 16t$
\end{tasks}

\begin{solution}
\textbf{Conceptual Explanation:} \\
"Ground level" means the initial height $h_0$ is 0. Substitute the given initial velocity for $v_0$.

\vspace{20pt}

\textbf{Calculation and Logic:} \\
$v_0 = 50$ \\
$h_0 = 0$ \\
$h(t) = -16t^2 + 50t + 0$

\vspace{25pt}

Answer: \colorbox{yellow!100}{A}
\end{solution}

\vspace{40pt}

% Question 20
% Category: Exam level | Difficulty: Hard | Question Type: Numeric Entry
\question
If the vertex of the parabola $y = x^2 + 8x + c$ is $(-4, 10)$, what is the value of $c$?

\begin{solution}
\textbf{Conceptual Explanation:} \\
Since the vertex is a point on the parabola, the coordinates must satisfy the equation. Substitute $x = -4$ and $y = 10$ to solve for $c$.

\vspace{20pt}

\textbf{Calculation and Logic:} \\
$10 = (-4)^2 + 8(-4) + c$ \\
$10 = 16 - 32 + c$ \\
$10 = -16 + c$ \\
$c = 26$.

\vspace{25pt}

Answer: \colorbox{yellow!100}{26}
\end{solution}

\vspace{40pt}

% Question 21
% Category: Practice | Difficulty: Medium | Question Type: Multiple Choice
\question
Which of the following functions has a graph with a vertex at $(3, -7)$?
\begin{tasks}(2)
    \task $y = (x-3)^2 - 7$
    \task $y = (x+3)^2 - 7$
    \task $y = (x-7)^2 + 3$
    \task $y = (x+7)^2 + 3$
\end{tasks}

\begin{solution}
\textbf{Conceptual Explanation:} \\
In vertex form $y = a(x-h)^2 + k$, the vertex is $(h, k)$.

\vspace{20pt}

\textbf{Calculation and Logic:} \\
For vertex $(3, -7)$: \\
$h = 3$ and $k = -7$. \\
$y = (x - 3)^2 - 7$.

\vspace{25pt}

Answer: \colorbox{yellow!100}{A}
\end{solution}

\vspace{40pt}

% Question 22
% Category: Concept | Difficulty: Medium | Question Type: Multiple Choice
\question
$f(x) = (x+4)(x-8)$ \\
At which point does the graph of the function above intersect the $y$-axis?
\begin{tasks}(4)
    \task $(0, -32)$
    \task $(0, 32)$
    \task $(-4, 0)$
    \task $(8, 0)$
\end{tasks}

\begin{solution}
\textbf{Conceptual Explanation:} \\
The $y$-intercept occurs where $x = 0$.

\vspace{20pt}

\textbf{Calculation and Logic:} \\
$f(0) = (0+4)(0-8)$ \\
$f(0) = 4(-8) = -32$.

\vspace{25pt}

Answer: \colorbox{yellow!100}{A}
\end{solution}

\vspace{40pt}

% Question 23
% Category: Exam level | Difficulty: Hard | Question Type: Numeric Entry
\question
$y = x^2 - kx + 9$ \\
The parabola above is tangent to the $x$-axis. If $k > 0$, what is the value of $k$?

\begin{solution}
\textbf{Conceptual Explanation:} \\
"Tangent to the $x$-axis" means there is exactly one solution, so the discriminant $b^2 - 4ac$ must be zero.

\vspace{20pt}

\textbf{Calculation and Logic:} \\
$a = 1, b = -k, c = 9$. \\
$(-k)^2 - 4(1)(9) = 0$ \\
$k^2 - 36 = 0$ \\
$k = \pm 6$.

\vspace{10pt}

Since $k > 0$, $k = 6$.

\vspace{25pt}

Answer: \colorbox{yellow!100}{6}
\end{solution}

\vspace{40pt}

% Question 24
% Category: Practice | Difficulty: Hard | Question Type: Multiple Choice
\question
\begin{center}
\begin{tikzpicture}[scale=0.6]
    \draw[thick] (0,0) -- (4,0) -- (4,3) -- (0,0);
    \node at (2,-0.4) {$x+2$};
    \node at (4.5, 1.5) {$x-1$};
\end{tikzpicture}
\end{center}
The area of the right triangle shown above is 10 square units. What is the value of $x$?
\begin{tasks}(4)
    \task 3
    \task 4
    \task 5
    \task 6
\end{tasks}

\begin{solution}
\textbf{Conceptual Explanation:} \\
Area of a triangle $= 0.5 \times Base \times Height$. Set up the equation and solve the quadratic.

\vspace{20pt}

\textbf{Calculation and Logic:} \\
$0.5(x+2)(x-1) = 10$ \\
$(x+2)(x-1) = 20$ \\
$x^2 + x - 2 = 20$

\vspace{10pt}

$x^2 + x - 22 = 0$. (Wait, checking for integer solutions). \\
Let's use Area $= 9$: \\
$0.5(x+2)(x-1) = 9 \Rightarrow x^2 + x - 2 = 18 \Rightarrow x^2 + x - 20 = 0 \Rightarrow (x+5)(x-4)=0$. \\
$x = 4$.

\vspace{25pt}

Answer: \colorbox{yellow!100}{B}
\end{solution}

\vspace{40pt}

% Question 25
% Category: Exam level | Difficulty: Hard | Question Type: Numeric Entry
\question
$f(x) = (x-3)^2 + 5$ \\
If $f(x) = f(k)$ and $k \neq x$, what is the sum of $x$ and $k$?

\begin{solution}
\textbf{Conceptual Explanation:} \\
In a parabola, two different $x$-values share the same $y$-value if they are symmetric about the axis of symmetry $x=h$. The average of these $x$-values is $h$.

\vspace{20pt}

\textbf{Calculation and Logic:} \\
Axis of symmetry: $h = 3$. \\
$\frac{x + k}{2} = 3$ \\
$x + k = 6$.

\vspace{25pt}

Answer: \colorbox{yellow!100}{6}
\end{solution}

% --- PART 2: Advanced Quadratic Word Problems (26-50) ---

% Question 26
% Category: Exam level | Difficulty: Hard | Question Type: Numeric Entry
\question
The revenue $R$, in dollars, for a small theater is modeled by the function $R(x) = -5(x - 40)^2 + 8000$, where $x$ is the price of a ticket in dollars. Based on the model, what ticket price, in dollars, results in the maximum revenue?

\begin{solution}
\textbf{Conceptual Explanation:} \\
In vertex form $y = a(x-h)^2 + k$, the maximum value (if $a < 0$) occurs at the $x$-coordinate of the vertex, which is $h$.

\vspace{20pt}

\textbf{Calculation and Logic:} \\
Identify the vertex $(h, k)$ from $R(x) = -5(x - 40)^2 + 8000$: \\
$h = 40$ \\
$k = 8000$

\vspace{10pt}

The variable $x$ represents the ticket price, and $R(x)$ represents revenue. The maximum revenue occurs when the ticket price is \$40.

\vspace{25pt}

Answer: \colorbox{yellow!100}{40}
\end{solution}

\vspace{40pt}

% Question 27
% Category: Practice | Difficulty: Medium | Question Type: Multiple Choice
\question
A rectangular picture has a length that is 4 inches greater than its width. If the width is $w$ inches, which of the following represents the area $A$ of the picture in square inches?
\begin{tasks}(4)
    \task $A = w^2 + 4$
    \task $A = w^2 + 4w$
    \task $A = 4w^2$
    \task $A = 2w + 4$
\end{tasks}

\begin{solution}
\textbf{Conceptual Explanation:} \\
Area is calculated as $Length \times Width$. We must express the length in terms of the width $w$ and multiply.

\vspace{20pt}

\textbf{Calculation and Logic:} \\
Width $= w$ \\
Length $= w + 4$ \\
Area $= w(w + 4) = w^2 + 4w$

\vspace{25pt}

Answer: \colorbox{yellow!100}{B}
\end{solution}

\vspace{40pt}

% Question 28
% Category: Concept | Difficulty: Medium | Question Type: Multiple Choice
\question
\begin{center}
\begin{tikzpicture}[scale=0.7]
    \draw[thick] (0,0) rectangle (5,3);
    \node at (2.5,-0.4) {$x+10$};
    \node at (-0.5, 1.5) {$x$};
\end{tikzpicture}
\end{center}
The area of the rectangle shown above is 75 square units. Which of the following equations can be used to find the value of $x$?
\begin{tasks}(2)
    \task $x^2 + 10 = 75$
    \task $2x + 10 = 75$
    \task $x^2 + 10x = 75$
    \task $x + (x+10) = 75$
\end{tasks}

\begin{solution}
\textbf{Conceptual Explanation:} \\
Setting the area product $x(x+10)$ equal to the given numerical value provides the equation for $x$.

\vspace{20pt}

\textbf{Calculation and Logic:} \\
Area $= Width \times Length$ \\
$75 = x(x+10)$ \\
$75 = x^2 + 10x$

\vspace{25pt}

Answer: \colorbox{yellow!100}{C}
\end{solution}

\vspace{40pt}

% Question 29
% Category: Exam level | Difficulty: Hard | Question Type: Numeric Entry
\question
An object is launched from a height of 10 feet with an initial upward velocity of 48 feet per second. The height $h$ is given by $h = -16t^2 + 48t + 10$. How many seconds after launch does the object reach its maximum height?

\begin{solution}
\textbf{Conceptual Explanation:} \\
The maximum height of a projectile occurs at the time $t$ corresponding to the $x$-coordinate of the vertex ($t = -b/2a$).

\vspace{20pt}

\textbf{Calculation and Logic:} \\
$a = -16, b = 48$ \\
$t = -48 / (2 \times -16)$ \\
$t = -48 / -32$ \\
$t = 1.5$

\vspace{25pt}

Answer: \colorbox{yellow!100}{1.5}
\end{solution}

\vspace{40pt}

% Question 30
% Category: Practice | Difficulty: Hard | Question Type: Multiple Choice
\question
\begin{center}
\begin{tikzpicture}[scale=0.6]
    \draw[->] (-1,0) -- (6,0) node[right] {$x$};
    \draw[->] (0,-1) -- (0,5) node[above] {$y$};
    \draw[domain=0.5:5.5, smooth, variable=\x, blue, thick] plot ({\x}, {-0.5*(\x-3)^2 + 4});
    \filldraw[black] (3,4) circle (2pt) node[above] {$(3, 4)$};
\end{tikzpicture}
\end{center}
The graph above models the path of a ball. If $y$ is the height in feet and $x$ is the horizontal distance in feet, what is the maximum height reached by the ball?
\begin{tasks}(4)
    \task 3 feet
    \task 4 feet
    \task 7 feet
    \task 12 feet
\end{tasks}

\begin{solution}
\textbf{Conceptual Explanation:} \\
The maximum height corresponds to the $y$-value of the vertex on the graph.

\vspace{20pt}

\textbf{Calculation and Logic:} \\
The vertex is labeled as $(3, 4)$. \\
$x = 3$ (horizontal distance). \\
$y = 4$ (maximum height).

\vspace{25pt}

Answer: \colorbox{yellow!100}{B}
\end{solution}

\vspace{40pt}

% Question 31
% Category: Exam level | Difficulty: Hard | Question Type: Numeric Entry
\question
The product of two consecutive positive even integers is 168. What is the value of the larger integer?

\begin{solution}
\textbf{Conceptual Explanation:} \\
Consecutive even integers differ by 2. Let the smaller be $x$ and the larger be $x+2$. Solve the quadratic equation formed by their product.

\vspace{20pt}

\textbf{Calculation and Logic:} \\
$x(x+2) = 168$ \\
$x^2 + 2x - 168 = 0$ \\
Factor: $(x+14)(x-12) = 0$ \\
$x = 12$ (since it must be positive).

\vspace{10pt}

Larger integer $= 12 + 2 = 14$.

\vspace{25pt}

Answer: \colorbox{yellow!100}{14}
\end{solution}

\vspace{40pt}

% Question 32
% Category: Concept | Difficulty: Medium | Question Type: Multiple Choice
\question
A projectile is fired upward. The function $h(t) = -16t^2 + 128t$ models its height. Which of the following represents the time $t$ when the projectile hits the ground?
\begin{tasks}(4)
    \task $t = 0$ only
    \task $t = 4$ only
    \task $t = 8$ only
    \task $t = 0$ and $t = 8$
\end{tasks}

\begin{solution}
\textbf{Conceptual Explanation:} \\
Hitting the ground means $h(t) = 0$. Factoring the quadratic will reveal the times at which the height is zero.

\vspace{20pt}

\textbf{Calculation and Logic:} \\
$-16t^2 + 128t = 0$ \\
$-16t(t - 8) = 0$ \\
$t = 0$ (at launch) and $t = 8$ (return to ground).

\vspace{25pt}

Answer: \colorbox{yellow!100}{D}
\end{solution}

\vspace{40pt}

% Question 33
% Category: Exam level | Difficulty: Hard | Question Type: Numeric Entry
\question
A rectangle has a perimeter of 40 centimeters. What is the maximum possible area, in square centimeters, of this rectangle?

\begin{solution}
\textbf{Conceptual Explanation:} \\
For a fixed perimeter, a rectangle reaches its maximum area when it is a square.

\vspace{20pt}

\textbf{Calculation and Logic:} \\
Perimeter $= 4s = 40$. \\
Side $s = 10$. \\
Area $= s^2 = 10^2 = 100$.

\vspace{25pt}

Answer: \colorbox{yellow!100}{100}
\end{solution}

\vspace{40pt}

% Question 34
% Category: Practice | Difficulty: Medium | Question Type: Multiple Choice
\question
$f(x) = (x-4)^2 - 9$ \\
Which of the following are the $x$-intercepts of the function $f$?
\begin{tasks}(4)
    \task $(4, 0)$ and $(-4, 0)$
    \task $(7, 0)$ and $(1, 0)$
    \task $(-7, 0)$ and $(-1, 0)$
    \task $(13, 0)$ and $(5, 0)$
\end{tasks}

\begin{solution}
\textbf{Conceptual Explanation:} \\
$x$-intercepts are found by setting $f(x) = 0$ and solving for $x$.

\vspace{20pt}

\textbf{Calculation and Logic:} \\
$(x-4)^2 - 9 = 0$ \\
$(x-4)^2 = 9$ \\
$x - 4 = \pm 3$ \\
$x = 4+3=7$ and $x = 4-3=1$.

\vspace{25pt}

Answer: \colorbox{yellow!100}{B}
\end{solution}

\vspace{40pt}

% Question 35
% Category: Concept | Difficulty: Medium | Question Type: Numeric Entry
\question
\begin{center}
\begin{tikzpicture}[scale=0.8]
    \draw[thick] (0,0) -- (6,0) -- (6,4) -- cycle;
    \node at (3,-0.4) {$x+4$};
    \node at (6.6, 2) {$x$};
\end{tikzpicture}
\end{center}
The area of the triangle above is 16. What is the value of $x$?

\begin{solution}
\textbf{Conceptual Explanation:} \\
Apply the area formula for a triangle: $Area = 0.5 \times Base \times Height$.

\vspace{20pt}

\textbf{Calculation and Logic:} \\
$16 = 0.5(x)(x+4)$ \\
$32 = x^2 + 4x$ \\
$x^2 + 4x - 32 = 0$ \\
$(x+8)(x-4) = 0$ \\
$x = 4$ (positive side length).

\vspace{25pt}

Answer: \colorbox{yellow!100}{4}
\end{solution}

\vspace{40pt}

% Question 36
% Category: Exam level | Difficulty: Hard | Question Type: Multiple Choice
\question
A company's daily cost $C$ to produce $x$ items is $C(x) = x^2 - 100x + 3000$. How many items should be produced to minimize the daily cost?
\begin{tasks}(4)
    \task 50
    \task 100
    \task 2,500
    \task 3,000
\end{tasks}

\begin{solution}
\textbf{Conceptual Explanation:} \\
The minimum of an upward-opening parabola occurs at the vertex $x = -b/2a$.

\vspace{20pt}

\textbf{Calculation and Logic:} \\
$a = 1, b = -100$ \\
$x = -(-100) / (2 \times 1) = 50$.

\vspace{25pt}

Answer: \colorbox{yellow!100}{A}
\end{solution}

\vspace{40pt}

% Question 37
% Category: Concept | Difficulty: Medium | Question Type: Numeric Entry
\question
If the graph of $y = (x-h)^2 + 5$ passes through $(2, 6)$, what is a possible value of $h$?

\begin{solution}
\textbf{Conceptual Explanation:} \\
Substitute the point into the equation and solve for the unknown constant $h$.

\vspace{20pt}

\textbf{Calculation and Logic:} \\
$6 = (2-h)^2 + 5$ \\
$1 = (2-h)^2$ \\
$\pm 1 = 2 - h$ \\
$h = 2-1=1$ or $h = 2+1=3$.

\vspace{25pt}

Answer: \colorbox{yellow!100}{1}
\end{solution}

\vspace{40pt}

% Question 38
% Category: Practice | Difficulty: Hard | Question Type: Multiple Choice
\question
Which function represents a parabola that opens downward and has a vertex at $(2, 10)$?
\begin{tasks}(2)
    \task $y = (x-2)^2 + 10$
    \task $y = -(x-2)^2 + 10$
    \task $y = -(x+2)^2 + 10$
    \task $y = -(x-10)^2 + 2$
\end{tasks}

\begin{solution}
\textbf{Conceptual Explanation:} \\
Vertex form is $a(x-h)^2+k$. "Downward" requires $a < 0$.

\vspace{20pt}

\textbf{Calculation and Logic:} \\
$h = 2, k = 10$. \\
$a$ must be negative. \\
$y = -(x-2)^2 + 10$.

\vspace{25pt}

Answer: \colorbox{yellow!100}{B}
\end{solution}

\vspace{40pt}

% Question 39
% Category: Concept | Difficulty: Medium | Question Type: Multiple Choice
\question
A ball is thrown from a height of 6 feet. Which of the following could be the equation for its height $h$ after $t$ seconds?
\begin{tasks}(2)
    \task $h = -16t^2 + 20t + 6$
    \task $h = -16t^2 + 6t + 20$
    \task $h = 6t^2 - 16t + 20$
    \task $h = -16t^2 + 20t - 6$
\end{tasks}

\begin{solution}
\textbf{Conceptual Explanation:} \\
The initial height is the constant term $c$ in standard form $y = at^2+bt+c$.

\vspace{20pt}

\textbf{Calculation and Logic:} \\
Initial height $= 6$. The constant term must be 6.

\vspace{25pt}

Answer: \colorbox{yellow!100}{A}
\end{solution}

\vspace{40pt}

% Question 40
% Category: Exam level | Difficulty: Hard | Question Type: Numeric Entry
\question
The sum of a number and its square is 72. What is the positive value of the number?

\begin{solution}
\textbf{Conceptual Explanation:} \\
Set up the quadratic $x + x^2 = 72$ and solve.

\vspace{20pt}

\textbf{Calculation and Logic:} \\
$x^2 + x - 72 = 0$ \\
$(x+9)(x-8) = 0$ \\
$x = 8$ (positive value).

\vspace{25pt}

Answer: \colorbox{yellow!100}{8}
\end{solution}

\vspace{40pt}

% Question 41
% Category: Practice | Difficulty: Hard | Question Type: Multiple Choice
\question
A square's side length is increased by 3. The new area is 100. What was the original side length?
\begin{tasks}(4)
    \task 7
    \task 10
    \task 13
    \task 49
\end{tasks}

\begin{solution}
\textbf{Conceptual Explanation:} \\
$(s+3)^2 = 100$.

\vspace{20pt}

\textbf{Calculation and Logic:} \\
$s + 3 = 10$ \\
$s = 7$.

\vspace{25pt}

Answer: \colorbox{yellow!100}{A}
\end{solution}

\vspace{40pt}

% Question 42
% Category: Exam level | Difficulty: Hard | Question Type: Numeric Entry
\question
What is the $y$-coordinate of the vertex of $y = x^2 - 6x + 14$?

\begin{solution}
\textbf{Conceptual Explanation:} \\
Find the vertex $x$, then plug it in to find the vertex $y$ (the minimum value).

\vspace{20pt}

\textbf{Calculation and Logic:} \\
$x = 6/2 = 3$. \\
$y = 3^2 - 6(3) + 14 = 9 - 18 + 14 = 5$.

\vspace{25pt}

Answer: \colorbox{yellow!100}{5}
\end{solution}

\vspace{40pt}

% Question 43
% Category: Concept | Difficulty: Medium | Question Type: Multiple Choice
\question
If a parabola is tangent to the $x$-axis, how many solutions does the equation $f(x) = 0$ have?
\begin{tasks}(4)
    \task Zero
    \task Exactly one
    \task Exactly two
    \task Infinitely many
\end{tasks}

\begin{solution}
\textbf{Conceptual Explanation:} \\
Tangency means the vertex touches the axis at exactly one point.

\vspace{20pt}

\textbf{Calculation and Logic:} \\
An $x$-intercept is a solution. Tangency implies one point of contact.

\vspace{25pt}

Answer: \colorbox{yellow!100}{B}
\end{solution}

\vspace{40pt}

% Question 44
% Category: Exam level | Difficulty: Hard | Question Type: Numeric Entry
\question
$f(x) = x^2 + kx + 16$ \\
If the vertex is on the $x$-axis and $k > 0$, what is the value of $k$?

\begin{solution}
\textbf{Conceptual Explanation:} \\
Vertex on $x$-axis means discriminant $b^2-4ac = 0$.

\vspace{20pt}

\textbf{Calculation and Logic:} \\
$k^2 - 4(1)(16) = 0$ \\
$k^2 = 64 \Rightarrow k = 8$.

\vspace{25pt}

Answer: \colorbox{yellow!100}{8}
\end{solution}

\vspace{40pt}

% Question 45
% Category: Practice | Difficulty: Medium | Question Type: Multiple Choice
\question
Which point is the vertex of $y = (x+5)^2 - 3$?
\begin{tasks}(4)
    \task $(5, -3)$
    \task $(-5, -3)$
    \task $(-5, 3)$
    \task $(0, -3)$
\end{tasks}

\begin{solution}
\textbf{Conceptual Explanation:} \\
Vertex is $(h, k)$ from $a(x-h)^2+k$.

\vspace{20pt}

\textbf{Calculation and Logic:} \\
$h = -5, k = -3$.

\vspace{25pt}

Answer: \colorbox{yellow!100}{B}
\end{solution}

\vspace{40pt}

% Question 46
% Category: Exam level | Difficulty: Hard | Question Type: Numeric Entry
\question
$y = -x^2 + 10x$ \\
What is the maximum value of this function?

\begin{solution}
\textbf{Conceptual Explanation:} \\
Find vertex $x$, then plug it in.

\vspace{20pt}

\textbf{Calculation and Logic:} \\
$x = -10 / -2 = 5$. \\
$y = -(5)^2 + 10(5) = -25 + 50 = 25$.

\vspace{25pt}

Answer: \colorbox{yellow!100}{25}
\end{solution}

\vspace{40pt}

% Question 47
% Category: Concept | Difficulty: Medium | Question Type: Multiple Choice
\question
A rectangle has length $L$ and width $W$. If $L + W = 10$, which function models the area $A$ in terms of $L$?
\begin{tasks}(2)
    \task $A = L(10-L)$
    \task $A = L(L-10)$
    \task $A = 10L$
    \task $A = L^2 + 10$
\end{tasks}

\begin{solution}
\textbf{Conceptual Explanation:} \\
Substitute $W = 10 - L$ into $A = LW$.

\vspace{25pt}

Answer: \colorbox{yellow!100}{A}
\end{solution}

\vspace{40pt}

% Question 48
% Category: Exam level | Difficulty: Hard | Question Type: Numeric Entry
\question
$y = (x-2)(x-10)$ \\
What is the $x$-coordinate of the vertex?

\begin{solution}
\textbf{Conceptual Explanation:} \\
Average of intercepts.

\vspace{20pt}

\textbf{Calculation and Logic:} \\
$(2 + 10) / 2 = 6$.

\vspace{25pt}

Answer: \colorbox{yellow!100}{6}
\end{solution}

\vspace{40pt}

% Question 49
% Category: Practice | Difficulty: Medium | Question Type: Numeric Entry
\question
If $f(x) = 2x^2 + 1$, what is $f(5)$?

\begin{solution}
\textbf{Calculation and Logic:} \\
$2(25) + 1 = 51$.

\vspace{25pt}

Answer: \colorbox{yellow!100}{51}
\end{solution}

\vspace{40pt}

% Question 50
% Category: Exam level | Difficulty: Hard | Question Type: Numeric Entry
\question
The area of a square is 121. If the side length is $x+1$, what is the value of $x$?

\begin{solution}
\textbf{Calculation and Logic:} \\
$(x+1)^2 = 121 \Rightarrow x+1 = 11 \Rightarrow x = 10$.

\vspace{25pt}

Answer: \colorbox{yellow!100}{10}
\end{solution}

% =============================================================
% Question End
% =============================================================

\end{questions}
\begin{center}
\rule{.5\textwidth}{1pt}
\end{center}
\end{document}
